\documentclass[11pt]{article}
\usepackage[margin=0.92in]{geometry}
\usepackage{amsmath,amssymb,amsthm,array,booktabs}
\usepackage[T1]{fontenc}
\usepackage{lmodern}
\usepackage[hidelinks]{hyperref}
\newtheorem{theorem}{Theorem}
\newtheorem{lemma}{Lemma}
\newcommand{\Q}{\mathbb Q}
\newcommand{\Gal}{\operatorname{Gal}}
\newcommand{\disc}{\operatorname{disc}}
\newcommand{\Res}{\operatorname{Res}}
\newcommand{\ord}{\operatorname{ord}}
\title{Full 5-torsion of the regular-pentagon quintic pencil}
\author{Alec Kriebel\\\small Independent researcher\\
\small\href{https://orcid.org/0009-0001-9320-500X}{ORCID: 0009-0001-9320-500X}}
\date{October 4, 2026}
\begin{document}
\maketitle
\begin{abstract}
We compute all twenty-five geometric fifth-torsion points of the genus-one
normalization of the regular-pentagon quintic pencil asked about in
McCallum's Question 17 in the AIM workshop problem list on rational and
integral points. For a fixed equation normalization over $K=\Q(\sqrt5)$,
we give a cubic model, both birational maps, a degree-ten polynomial and
both ordinate branches for the twenty points outside the known infinity
subgroup. The complete division field is
$K(\sqrt{5+2\sqrt5},\zeta_5,\sqrt[5]{\lambda+\phi^5})$,
where $\phi=(1+\sqrt5)/2$. Its degree is four or twenty, with the exact
splitting criterion and Galois action specified. The formulas include
every elliptic normalization, including the allowed cuspidal plane
member. The method uses classical Tate normal form, division-polynomial
theory and Fisher's full-level cover. A reproducible verification
supplement accompanies this unrefereed, extensively AI-assisted note.
\end{abstract}

\section{The pencil and the answer}
The original question concerns the normalization of a pencil formed by
the five side lines of a regular pentagon and the square of its
circumcircle equation \cite[Question 17, p.~51]{aim}. The source already
observes the five-point infinity subgroup. We compute the other twenty
points and the extension data as well. We fix coordinates, scaling and
origin explicitly; changing the scaling changes the numerical pencil
parameter. We make no first-priority claim. The universal modular and
division-polynomial ingredients below are classical.

Put
\[
r=\sqrt5,\quad K=\Q(r),\quad \phi=(1+r)/2,\quad
c=\phi^5=(11+5r)/2,\quad d=5+2r,\quad \delta^2=d.
\]
Let $Q=X^2+Y^2-T^2$ and
\begin{equation}\label{pentagon}
P=2(X^5-10X^3Y^2+5XY^4)+5\phi(X^2+Y^2)^2T
 -5\phi^3(X^2+Y^2)T^3+cT^5.
\end{equation}
The projective pencil is $\Gamma_\lambda:P+\lambda TQ^2=0$.
The factor $T$ is the homogenization of the affine degree-four term.
Let $E_\lambda$ denote its smooth projective normalization, with origin
$O=[0:1:0]$.

\begin{theorem}\label{main}
In characteristic zero the finite members have a geometrically integral
genus-one normalization exactly when
\begin{equation}\label{allowed}
\lambda\notin\{0,-5r,-c\}.
\end{equation}
For every such member, equations \eqref{tate-map} and
\eqref{residual-table}--\eqref{ordinates}, together with the four points
\eqref{marked} and the origin, compute all $25$ geometric fifth-torsion
points. For $\lambda\in K$ satisfying \eqref{allowed}, choose a primitive
fifth root $\zeta$ and $\theta$ with $\theta^5=\lambda+c$. Then
\begin{equation}\label{field}
K(E_\lambda[5])=K(\delta,\zeta,\theta).
\end{equation}
This has degree $4$ when $\lambda+c\in K^{\times5}$ and degree $20$
otherwise. The respective Galois groups are $C_2\times C_2$ and
$D_{10}\times C_2$, where $D_{10}$ has order ten. Moreover
\begin{equation}\label{rational}
E_\lambda(K)[5]=0,\qquad
E_\lambda(K(\delta))[5]\cong\mathbb Z/5\mathbb Z,
\end{equation}
the latter subgroup being the five infinity points. Generically,
\eqref{field} holds over $K(\lambda)$ and has degree twenty.
\end{theorem}

The question's remarks about twists, local conditions and nonregular or
star polygons are separate variants. We do not assert a
Tate--Shafarevich construction or solve those variants. This specified
model has the rational origin $O$.

\section{Normalization and the exact elliptic range}
To check that \eqref{pentagon} is the side product, set $Z=X+iY$,
$W=X-iY$ and choose $\zeta=e^{2\pi i/5}$. The factors
$Z+\zeta^{2j+1}W-(1+\zeta)\zeta^jT$, for $0\le j<5$, vanish on
successive unit-circle vertices. Their product is
\[
Z^5+W^5+5\phi Z^2W^2T-5\phi^3ZWT^3+cT^5=P.
\]
Thus this is the required regular pentagon, with a fixed nonzero scaling.

In the chart $T=1$, write the equation as $AY^4+BY^2+C=0$, where
\begin{align*}
A&=10X+5\phi+\lambda,\\
B&=-20X^3+2(5\phi+\lambda)X^2-5\phi^3-2\lambda,\\
C&=2X^5+(5\phi+\lambda)X^4-(5\phi^3+2\lambda)X^2+c+\lambda.
\end{align*}
With $h=4X^2+2X-1$, $a=40+20r+8\lambda$, $b=a+5$, direct expansion gives
\begin{equation}\label{conic}
B^2-4AC=h^2(20X^2-aX+b),\quad
V=\frac{2AY^2+B}{h},\quad t=V-2rX,\quad
X=\frac{b-t^2}{a+4rt}.
\end{equation}
Put $\alpha=1-r/5$, $\gamma=2r/5$, $m=5-2r$, $z=t+d$, $s=z+\alpha\lambda$,
and $w=20Y(z+\gamma\lambda)/z$. Define
\begin{align*}
F(z)={}&z^3-((3+r)\lambda+40+20r)z^2\\
 &+((60+36r)\lambda+900+400r)z
 +(48+16r)\lambda^2+(500+220r)\lambda.
\end{align*}
Substitution in \eqref{conic} gives
\begin{equation}\label{double-cover}
Y^2=\frac{mz^2F(z)}{400(z+\gamma\lambda)^2(z+\alpha\lambda)},
\qquad w^2=\frac{mF(s-\alpha\lambda)}s.
\end{equation}
Set
\begin{align}
a_0&=5-2r,& a_1&=(-26+10r)\lambda-20r,\notag\\
a_2&=(42-86r/5)\lambda^2+(-100+100r)\lambda+500+200r,\notag\\
a_3&=(-112+48r)\lambda^2(\lambda+5r)/5.\label{coefficients}
\end{align}
The coordinates $\xi=a_3/s$, $\eta=\xi w$ yield the monic cubic
\begin{equation}\label{cubic}
W_\lambda:\quad \eta^2=\xi^3+a_2\xi^2+a_1a_3\xi+a_0a_3^2.
\end{equation}
The inverse map is
\begin{equation}\label{inverse}
s=a_3/\xi,\quad z=s-\alpha\lambda,\quad t=z-d,\quad
X=\frac{b-t^2}{a+4rt},\quad
Y=\frac{\eta z}{20\xi(z+\gamma\lambda)}.
\end{equation}
These are inverse function-field maps. For example, with
$G=AY^4+BY^2+C$ and the forward value of $t$, the exact identity
\[
b-t^2-X(a+4rt)=-4AG/h^2
\]
proves the recovered $X$ is unchanged; $Y$ cancels literally. Conversely
\eqref{cubic} recovers the conic branch in \eqref{double-cover} and hence
the same $t,\xi,\eta$. These fractions specify dense open charts; at
poles their unique extensions between smooth projective normalizations
are understood.

Here are the identities establishing the global, rather than selected
component, normalization:
\begin{align}
a^2-80b&=64\lambda(\lambda+5r),\notag\\
F(-\alpha\lambda)&=(-16/5+16r/25)\lambda^2(\lambda+5r),\notag\\
\disc(F)&=(24064+10752r)\lambda(\lambda+5r)^3(\lambda+c),\notag\\
\Delta_W&=2^{24}(161-72r)\lambda^5(\lambda+5r)^5(\lambda+c).
\label{discriminants}
\end{align}
Over an algebraic closure, the conic square root in \eqref{conic} is
nontrivial for $\lambda(\lambda+5r)\ne0$. Its field is rational in $t$.
Under \eqref{allowed}, the remaining square root in
\eqref{double-cover} has precisely four odd branch points: the three
distinct roots of $F$ and $s=0$. Thus the original quartic in $Y$ is
irreducible over the algebraic closure of $K$, and Riemann--Hurwitz gives
genus one. This also proves that \eqref{cubic} normalizes the whole
plane curve.

At $\lambda=0$ the curve is five side lines. At $\lambda=-5r$,
\eqref{pentagon} is the same line-product expression with $\phi$
replaced by $(1-r)/2$, again five lines. At $\lambda=-c$ the conic is
nonsingular and $F$ has exactly one double root and no triple root; the
gcd in the $s$-coordinate is $s-1-3r/5$. Removing its square leaves two
odd branch points, giving genus zero. The infinite pencil member is
$TQ^2=0$. By contrast, at $\lambda=-(25+10r)/4$ the five plane vertex
nodes become cusps but \eqref{discriminants} is nonzero, so this member
is included in the elliptic range.

At infinity on \eqref{cubic}, $\ord_O(\xi)=-2$, $\ord_O(\eta)=-3$,
so $s$ vanishes to order two and $w$ has a simple pole. In
\eqref{inverse}, $z=-\alpha\lambda$, $z+\gamma\lambda=(\gamma-\alpha)\lambda$
and $a+4rt=4(3-r)\lambda$ are nonzero. Therefore $Y$ has a simple pole
and $X=-(5\phi+\lambda)/10$ is finite: the image is $[0:1:0]$.
Its plane $X$ partial is ten, so it is smooth. Thus the origins agree.

\section{The actual Tate twist and all the torsion points}
Define
\begin{equation}\label{parameters}
\beta=\frac{(11-5r)\lambda}{2(\lambda+5r)},\qquad
k=\frac{4(\lambda+5r)}r,\qquad q=dk^2.
\end{equation}
The Tate equation \cite[Lemma 1.1 and (2), p.~172]{fisher} is
\begin{equation}\label{tate}
D_\beta:\quad y^2+(1-\beta)xy-\beta y=x^3-\beta x^2,
\qquad \Delta_\beta=\beta^5(\beta^2-11\beta-1).
\end{equation}
The following is an origin-preserving equation isomorphism over
$K(\delta)$, valid on every allowed fiber:
\begin{equation}\label{tate-map}
\xi=q(x-\beta),\qquad
\eta=(k\delta)^3\left(y+\frac{(1-\beta)x-\beta}{2}\right).
\end{equation}
Indeed, put
\begin{equation}\label{branch}
T_\beta(x)=4x^3+(\beta^2-6\beta+1)x^2+2(\beta^2-\beta)x+\beta^2.
\end{equation}
The right-hand side of \eqref{cubic} at $\xi=q(x-\beta)$ is
$q^3T_\beta(x)/4$, and $(k\delta)^6=q^3$.
This verifies the quadratic twist, including its sign, rather than just
a $j$-invariant equality. No extra exclusion at $j=0$ or $1728$ occurs.

The known nonzero points on $D_\beta$ are
\begin{equation}\label{marked}
(0,0),\quad(0,\beta),\quad(\beta,0),\quad(\beta,\beta^2).
\end{equation}
For $P_0=(0,0)$ the tangent is horizontal; the chord law gives
$2P_0=(\beta,\beta^2)$ and $3P_0=(\beta,0)=-2P_0$, proving exact order
five. The other points are its nonzero multiples.

Let $R_\beta(x)=\sum_{j=0}^{10}r_jx^j$ with the following coefficients:
\begin{equation}\label{residual-table}
\begin{array}{c|l}
j&r_j\\\hline
10&5\\
9&5(\beta^2-5\beta+1)\\
8&\beta^4-7\beta^3+44\beta^2-38\beta+1\\
7&\beta(\beta^4+3\beta^3-26\beta^2+127\beta-9)\\
6&\beta^2(\beta^4+3\beta^3+19\beta^2-248\beta+36)\\
5&\beta^3(\beta^4+3\beta^3-71\beta^2+322\beta-84)\\
4&\beta^4(\beta^4-12\beta^3+94\beta^2-293\beta+126)\\
3&5\beta^5(\beta^3-10\beta^2+36\beta-25)\\
2&5\beta^6(2\beta^2-13\beta+16)\\
1&10\beta^7(\beta-3)\\
0&5\beta^8.
\end{array}
\end{equation}
For every root $x$ take both ordinates
\begin{equation}\label{ordinates}
y=\frac{-((1-\beta)x-\beta)\pm\sqrt{T_\beta(x)}}2.
\end{equation}
Then apply \eqref{tate-map} and, if desired, \eqref{inverse}.

\begin{lemma}\label{complete}
The ten roots of $R_\beta$ and the two choices in \eqref{ordinates}
are distinct on every allowed fiber and give exactly the twenty
fifth-torsion points outside \eqref{marked} and the origin.
\end{lemma}
\begin{proof}
Put $b_2=\beta^2-6\beta+1$, $b_4=\beta^2-\beta$, $b_6=\beta^2$,
$b_8=-\beta^3$. The classical division polynomials give
\begin{align*}
\psi_3&=3x^4+b_2x^3+3b_4x^2+3b_6x+b_8,\\
H_6&=2x^6+b_2x^5+5b_4x^4+10b_6x^3+10b_8x^2
 +(b_2b_8-b_4b_6)x+b_4b_8-b_6^2,\\
\psi_5&=T_\beta^2H_6-\psi_3^3=x(x-\beta)R_\beta.
\end{align*}
These identities can also be derived directly from the chord law on
$v^2=f(x)$: writing $x_2=x(2P)$, $x_3=x(3P)$ gives
$\psi_3=4f(x-x_2)$, $H_6=16f^2v(2P)/v(P)$ and
$\psi_5=4f\psi_3^2(x_2-x_3)$. The supplement performs that independent
generic calculation and checks each coefficient in \eqref{residual-table}.
The exact resultants and evaluations are
\[
\Res_x(\psi_5,T_\beta)=\Delta_\beta^6,\quad
\Res_x(\psi_5,\psi_3)=\Delta_\beta^8,\quad
R_\beta(0)=5\beta^8,\quad R_\beta(\beta)=5\beta^{12}.
\]
For a finite point with $\psi_5=0$ these nonzero resultants justify
the chord denominators and imply $x(2P)=x(3P)$. Equality $2P=3P$
would imply $P=O$, impossible, so $3P=-2P$ and $5P=O$.
Conversely a nonzero fifth-torsion point has those denominators nonzero
and satisfies this equation. In characteristic zero $[5]$ is separable
of degree 25 \cite[sections 5.5--5.6]{sutherland}; hence its kernel has
25 distinct points. Negation pairs the24 nonzero points into twelve
distinct abscissae. Thus $\psi_5$ has twelve simple roots, and the
resultant with $T_\beta$ makes both ordinate signs distinct. Removing
the two marked abscissae leaves exactly the asserted twenty points.
\end{proof}

There is no additional inverse-chart exclusion for these twenty points.
After \eqref{tate-map}, direct substitution gives
\[
z+\gamma\lambda=(\gamma-\alpha)\lambda\frac{x}{x-\beta},
\qquad a+4rt=4r(z+\gamma\lambda).
\]
Thus the poles $\xi=0$ and $z+\gamma\lambda=0$ occur only at the two
known abscissae, which are not roots of $R_\beta$. The five infinity
points themselves have $X/Y=0,\pm\sqrt{5+2r},\pm\sqrt{5-2r}$.
They are smooth, and rotation by 72 degrees preserves the pencil and
cycles them. Its action on the normalization is an order-five
translation, since the origin-fixing automorphism group in
characteristic zero has no element of order five. They therefore form
the order-five subgroup in \eqref{rational}, in agreement with the
source's infinity observation.

\section{The full division field and specialization}
We specify the modular input; the infinity line or a semisimplification
alone would not determine this field. Choose $\zeta$ with
$1+\zeta+\zeta^{-1}=\phi$. Over a field containing $\zeta$, the Tate
curve with marked $P_0$ has five choices of $Q_0$ satisfying
$e_5(P_0,Q_0)=\zeta$, namely $Q_0+jP_0$. These form the finite etale
degree-five full-level cover of the marked Tate parameter line away
from its cusps. An elliptic automorphism fixing a point of order at
least four is trivial \cite[Lemma 1.1]{fisher}; thus the cover is fine.
Fisher's full-level description \cite[section 2 and Lemma 3.4,
p.~194]{fisher} gives its parameter map
\[
\beta=\tau\frac{f(\tau)}{g(\tau)},\quad
f=\tau^4+3\tau^3+4\tau^2+2\tau+1,\quad
g=\tau^4-2\tau^3+4\tau^2-3\tau+1.
\]
With $\varepsilon(\tau)=(\phi\tau+1)/(\tau-\phi)$ and
$\iota(v)=(cv+1)/(v-c)$, both transformations are involutions and
the exact identities are
\begin{equation}\label{kummer}
\tau f(\tau)/g(\tau)=\iota(\varepsilon(\tau)^5),\qquad
\iota(\beta)=-\frac1{\lambda+c}.
\end{equation}
Consequently the full-level cover is $u^5=-1/(\lambda+c)$, with
$u=\varepsilon(\tau)$. This identifies every allowed specialization,
not merely a generic extension: both covers are the normalization of
the same function-field cover over the smooth cusp-complement base,
and both are finite etale there. The fiber is exactly the choices of
the complementary basis vector, including split fibers and the
unmarked exceptional $j$ values. No generic irreducibility assertion
is imposed on every specialization.

Write $a=\lambda+c\ne0$. If $\theta^5=a$, then $u=-1/\theta$ satisfies
$u^5=-1/a$. Thus the two radical fields coincide. For a fixed
identification of the Kummer group, the classes satisfy
$[-1/a]=[a]^{-1}$, rather than being equal as nontrivial fixed classes;
the inversion of generator or coordinate must be stated.

Let $L=K(\delta,\zeta)$. The fine-cover interpretation and
\eqref{tate-map} give $L(E_\lambda[5])=L(\theta)$ in both directions:
a cover point supplies a basis and a basis determines its cover point.
The image of $P_0$ has
\[
\xi_{P_0}=-q\beta\in K,\qquad
\eta_{P_0}=-k^3d\delta\beta/2.
\]
The latter has a nonzero $K$ coefficient, so the full torsion field
contains $\delta$. Nondegeneracy of the Weil pairing supplies
$\zeta$ \cite[Theorem 23.29]{pairing}. Hence it already contains $L$,
proving \eqref{field}.

For completeness, $M=K(\delta)$ is quadratic because the norm of $d$
to $\Q$ is five, not a rational square. It is real, while
$K(\zeta)/K$ is imaginary quadratic; hence $[L:K]=4$.
If $v\in L$ and $v^5=a\in K^*$, then
$N_{L/K}(v)^5=a^4$, so $(a/N_{L/K}(v))^5=a$. Thus $a$ becomes a
fifth power in $L$ exactly when it is one in $K$. Kummer theory now
gives degrees 4 or 20. In the nonsplit case $K(\zeta,\theta)/K$ has
group $D_{10}$: a rotation multiplies $\theta$ by $\zeta$ and the
involution fixing the real fifth root inverts $\zeta$. Its unique
quadratic subfield is $K(\zeta)$, distinct from $M$, so adjoining
$\delta$ gives $D_{10}\times C_2$. The same proof works generically;
the valuation of $\lambda+c$ at $\lambda=-c$ is one, so it is not a
fifth power in $K(\lambda)$.

The independent involution changing $\delta$ and fixing the Tate
torsion field acts by $-I$ under \eqref{tate-map}. Therefore there is
no nonzero $K$-rational fifth torsion. Over the real field $M$ the
known subgroup exists and $E(\mathbb R)[5]$ has exactly five elements,
so it is the entire $M$-rational subgroup. Finally, choosing the
infinity generator as the first basis vector and adjusting the second
to be anti-invariant under complex conjugation gives the matrices
\[
\sigma_\delta\mapsto-I,\qquad
\sigma_\zeta\mapsto\begin{pmatrix}1&0\\0&-1\end{pmatrix},\qquad
\sigma_5\mapsto\begin{pmatrix}1&1\\0&1\end{pmatrix}.
\]
The nonzero transvection is normalized by choosing its generator. In
the split degree-four case it is omitted. These matrices give the
full extension, with the basis and generator choices made explicit.

\section{Verification and research status}
The accompanying package supplies exact symbolic geometry and
generic chord-law checks, independent finite-field group-law
enumerations, the arithmetic field controls, and deliberate mutants.
Universal proofs establish the characteristic-zero all-fiber claims;
finite tests are falsifiable implementation checks, not proofs by
sampling. The arithmetic controls are also reproducible with the
Python standard library. The package records dependencies, expected
mathematical outputs and source references without redistributing
third-party primary PDFs.

This note assembles an explicit computation for the source's
regular-pentagon pencil, extending the already observed infinity
subgroup to the complete kernel and division field. Tate normal form,
Fisher's modular description, the division-polynomial theory and the
Weil pairing are credited above. The bounded priority audit in the
supplement documents its actual searches and reading limits; no
absence-of-hits argument certifies first discovery or present-day
openness. The repository's earlier partial infinity and quotient
calculation is acknowledged as an input to the original derivation.

\paragraph{AI use and review.}
AI tools were used extensively in derivation, computation, literature
research, writing, reproduction and adversarial review. The author
assumes responsibility for the submitted claims and artifacts. This
is an unrefereed preprint. Automated adversarial reviews do not
constitute independent external human peer review, and no such human
review or formal proof-assistant certification is claimed.

\begin{thebibliography}{9}
\bibitem{aim} American Institute of Mathematics,
\emph{Rational and integral points on higher dimensional varieties},
workshop problem list, December 11--20, 2002; version stamped
November 22, 2004, Question 17 and remarks, p.~51.
\url{https://www.aimath.org/WWN/qptsurface2/qptsurface2.pdf}.
\bibitem{fisher} T. A. Fisher, \emph{Some examples of 5 and 7 descent for
elliptic curves over $\Q$}, J. Eur. Math. Soc. \textbf{3} (2001), 169--201.
\href{https://doi.org/10.1007/s100970100030}{doi:10.1007/s100970100030}.
\bibitem{sutherland} A. V. Sutherland, \emph{18.783 Elliptic Curves,
Lecture 5: Isogeny kernels and division polynomials}, MIT, September 26, 2023,
sections 5.5--5.6.
\url{https://math.mit.edu/classes/18.783/2023/LectureNotes5.pdf}.
\bibitem{pairing} A. V. Sutherland, \emph{18.783 Elliptic Curves,
Lecture 23: Divisors and the Weil pairing}, MIT, December 5, 2023, section 23.5.
\url{https://math.mit.edu/classes/18.783/2023/LectureNotes23.pdf}.
\end{thebibliography}
\end{document}
